\section*{Capítulo \ref{ch:g10:muscles-and-joints} --- Músculos e articulações}

\begin{solution}{exo:g10:muscles-and-joints:1}
\omterm{def:g10:muscles-and-joints:joint}{Cartilagem} nas extremidades dos ossos (deslizamento suave, distribuição
da carga); cápsula (veda a \omterm{def:g10:muscles-and-joints:joint}{articulação}) com o líquido sinovial dela
(lubrificação); \omterm{def:g10:muscles-and-joints:joint}{ligamentos} (seguram os ossos juntos, limitam os
movimentos); no joelho, os meniscos (distribuem a carga).
\end{solution}

\begin{solution}{exo:g10:muscles-and-joints:2}
Um \omterm{def:g10:muscles-and-joints:muscle}{tendão} une um músculo a um osso e transmite a tração dele; um
\omterm{def:g10:muscles-and-joints:joint}{ligamento} une osso a osso através de uma \omterm{def:g10:muscles-and-joints:joint}{articulação} e limita o
movimento dela. Nenhum dos dois se contrai.
\end{solution}

\begin{solution}{exo:g10:muscles-and-joints:3}
Um músculo só consegue puxar encurtando; ele não consegue empurrar.
Para mover uma \omterm{def:g10:muscles-and-joints:joint}{articulação} nos dois sentidos são necessários um flexor
e um extensor, cada um desfazendo o que o outro fez.
\end{solution}

\begin{solution}{exo:g10:muscles-and-joints:4}
Entorse: um \omterm{def:g10:muscles-and-joints:joint}{ligamento} esticado ou rompido. Ruptura muscular: \omterm{def:g10:muscles-and-joints:muscle}{fibras
musculares} rompidas. Luxação: a extremidade de um osso forçada para
fora da \omterm{def:g10:muscles-and-joints:joint}{articulação}, com a cápsula e os \omterm{def:g10:muscles-and-joints:joint}{ligamentos} esticados ou
rompidos.
\end{solution}

\begin{solution}{exo:g10:muscles-and-joints:5}
Fibras lentas: lentas, fracas, incansáveis, ricas em \omterm{def:g10:cells-common-unit:organelle}{mitocôndrias} e em
capilares, avermelhadas, funcionam com a respiração. Fibras rápidas:
rápidas, potentes, cansam em um minuto, poucas \omterm{def:g10:cells-common-unit:organelle}{mitocôndrias}, pálidas,
funcionam com a reserva de fosfato e com a \omterm{def:g10:cell-metabolism:fermentation}{fermentação}.
\end{solution}

\begin{solution}{exo:g10:muscles-and-joints:6}
Carga $2 \times 9.8 = \qty{19.6}{N}$; força do bíceps $19.6 \times 32/4
\approx \qty{157}{N}$.
\end{solution}

\begin{solution}{exo:g10:muscles-and-joints:7}
Um empurrão pode vir de qualquer um dos dois lados, e cada músculo só
consegue resistir a um. Os dois contraídos enrijecem a \omterm{def:g10:muscles-and-joints:joint}{articulação} nos
dois sentidos ao mesmo tempo; de onde quer que venha o empurrão, um
deles já está puxando contra ele.
\end{solution}

\begin{solution}{exo:g10:muscles-and-joints:8}
A fibra encurta \qty{0.6}{cm}; o \omterm{def:g10:muscles-and-joints:muscle}{tendão} move o osso em \qty{0.6}{cm} a
\qty{5}{cm} do cotovelo, de modo que a mão se desloca
$0.6 \times 35/5 = \qty{4.2}{cm}$.
\end{solution}

\begin{solution}{exo:g10:muscles-and-joints:9}
A mobilidade exige uma cavidade rasa e \omterm{def:g10:muscles-and-joints:joint}{ligamentos} frouxos; essa mesma
frouxidão deixa a cabeça do úmero sair da cavidade quando uma força a
empurra além da amplitude que os \omterm{def:g10:muscles-and-joints:joint}{ligamentos} conseguem segurar.
\end{solution}

\begin{solution}{exo:g10:muscles-and-joints:10}
A desidratação e a perda de sal pela transpiração, e a fadiga das
fibras no fim da prova. Remédio: alongar o músculo com suavidade e
beber água com sal.
\end{solution}

\begin{solution}{exo:g10:muscles-and-joints:11}
A coxa do velocista tem 70\% de fibras rápidas e a proporção é em
grande parte herdada; o treinamento aumenta as fibras e melhora o
equipamento delas, mas não converte o número delas. A pessoa não
treinada parte de metade e metade e tem mais fibras lentas a
desenvolver.
\end{solution}

\begin{solution}{exo:g10:muscles-and-joints:12}
$4 \times 70 \times 9.8 \approx \qty{2.7}{kN}$. Tensão $2744/100
\approx \qty{27}{N/mm^2}$: cerca de um quarto da tensão de ruptura, um
fator de segurança de quatro.
\end{solution}

\begin{solution}{exo:g10:muscles-and-joints:13}
O músculo é ricamente irrigado de sangue e contém \omterm{def:g10:cells-common-unit:cell}{células} de reserva
que reconstroem fibras: semanas. Os \omterm{def:g10:muscles-and-joints:joint}{ligamentos} têm poucos vasos, de
modo que os nutrientes e as \omterm{def:g10:cells-common-unit:cell}{células} de reparo chegam devagar: meses. A
\omterm{def:g10:muscles-and-joints:joint}{cartilagem} não tem vaso sanguíneo nenhum e quase nenhuma \omterm{def:g10:cells-common-unit:cell}{célula} em
divisão; ela é alimentada só pelo líquido sinovial e quase não se
repara.
\end{solution}

\begin{solution}{exo:g10:muscles-and-joints:14}
A musculação engrossa as fibras (mais miofibrilas em cada uma) e o
alongamento aumenta o comprimento do conjunto músculo--tendão e a
tolerância da cápsula. O número de fibras é fixado no início da vida; o
músculo adulto cresce ou encolhe mudando o tamanho das fibras
existentes.
\end{solution}

\begin{solution}{exo:g10:muscles-and-joints:15}
Nas escadas a \omterm{def:g10:muscles-and-joints:joint}{articulação} carrega várias vezes o peso do corpo; uma
\omterm{def:g10:muscles-and-joints:joint}{cartilagem} mais fina transmite essa pressão ao osso de baixo, que tem
nervos e dói. Com menos \omterm{def:g10:muscles-and-joints:joint}{cartilagem}, os meniscos são a última coisa que
distribui a carga. Músculos fortes em volta do joelho absorvem parte do
impacto e estabilizam a \omterm{def:g10:muscles-and-joints:joint}{articulação}, reduzindo a pressão de pico sobre
a \omterm{def:g10:muscles-and-joints:joint}{cartilagem} que resta.
\end{solution}

\begin{solution}{pb:g10:muscles-and-joints:1}
\textbf{1.} $70 \times 9.8 = \qty{686}{N}$.

\textbf{2.} $686 \times 10/5 \approx \qty{1.4}{kN}$.

\textbf{3.} $686 \times 20/5 \approx \qty{2.7}{kN}$.

\textbf{4.} Duas e quatro vezes o peso do corpo: a cada degrau o
quadríceps puxa com o peso de quatro pessoas, e é por isso que as coxas
ardem numa subida longa.

\textbf{5.} Os isquiotibiais, atrás da coxa; eles relaxam e são
esticados enquanto o quadríceps estica o joelho, contraindo-se de leve
para estabilizar a \omterm{def:g10:muscles-and-joints:joint}{articulação}.

\textbf{6.} Força do chão $3 \times 686 \approx \qty{2.06}{kN}$; força
no \omterm{def:g10:muscles-and-joints:muscle}{tendão} $2058 \times 20/5 \approx \qty{8.2}{kN}$, doze vezes o peso
do corpo.

\textbf{7.} $8232/120 \approx \qty{69}{N/mm^2}$: cerca de dois terços
da tensão de ruptura. Cada aterrissagem usa quase toda a margem do
\omterm{def:g10:muscles-and-joints:muscle}{tendão}.

\textbf{8.} Carga $8.2 + 2.1 \approx \qty{10.3}{kN}$. Com os meniscos:
$10300/600 \approx \qty{17}{N/mm^2}$; sem eles: $10300/100 \approx
\qty{103}{N/mm^2}$, seis vezes mais.

\textbf{9.} Seis vezes a pressão sobre a mesma \omterm{def:g10:muscles-and-joints:joint}{cartilagem} a cada
aterrissagem: a superfície é esmagada mais depressa do que consegue ser
mantida, e se desgasta.

\textbf{10.} Uma película de líquido entre as cartilagens, para que
elas deslizem durante a flexão sem atrito seco, e um leve amortecimento
do impacto.

\textbf{11.} Uma torção da tíbia sob o fêmur, com uma flexão lateral:
rotações para as quais o joelho não é construído, e às quais os
\omterm{def:g10:muscles-and-joints:joint}{ligamentos} estão ali para resistir.

\textbf{12.} O quadríceps só puxa ao longo da perna, esticando o
joelho; ele não tem ação nenhuma contra uma torção. E a lesão acontece
em poucos centésimos de segundo, mais depressa que qualquer resposta
muscular.

\textbf{13.} Sangue dos vasos do \omterm{def:g10:muscles-and-joints:joint}{ligamento} rompido e da cápsula, mais
líquido sinovial em excesso: a cápsula foi rompida ou inflamou e a
\omterm{def:g10:muscles-and-joints:joint}{articulação} se encheu.

\textbf{14.} O \omterm{def:g10:muscles-and-joints:joint}{ligamento}, mal irrigado de sangue, precisa de meses e
muitas vezes de cirurgia; o menisco lesionado, \omterm{def:g10:muscles-and-joints:joint}{cartilagem} quase sem
vasos, em geral não cicatriza e a parte rompida é aparada ou suturada.

\textbf{15.} Os isquiotibiais são os antagonistas do quadríceps e puxam
a tíbia para trás; um par forte contraindo-se junto enrijece o joelho
contra o deslizamento para a frente que o \omterm{def:g10:muscles-and-joints:joint}{ligamento} rompido não impede
mais, protegendo a \omterm{def:g10:muscles-and-joints:joint}{articulação} antes e depois da cirurgia.

\textbf{16.} Cerca de 4000 vezes por dia a \qty{1.4}{kN} ou mais.

\textbf{17.} Cerca de $\num{1.5e6}$ solicitações por ano. A tendinite é
um dano que se acumula mais depressa que o reparo lento do tecido do
\omterm{def:g10:muscles-and-joints:muscle}{tendão}; só o descanso deixa o reparo alcançá-lo, já que a força não
reduz a carga de cada passo.

\textbf{18.} Aterrissagens $40 \times 8.2 \approx \qty{330}{kN}$ ao
todo; passos $4000 \times 1.4 \approx \qty{5600}{kN}$ --- os passos
carregam dezessete vezes mais carga total, mas as aterrissagens carregam
os picos que ameaçam romper.

\textbf{19.} Dobrar a força do chão dobra as forças no \omterm{def:g10:muscles-and-joints:muscle}{tendão} e na
\omterm{def:g10:muscles-and-joints:joint}{articulação} a cada passada; a bicicleta e a natação trabalham os
músculos carregando o joelho bem abaixo do peso do corpo.

\textbf{20.} O \omterm{def:g10:muscles-and-joints:muscle}{tendão} do quadríceps carrega quatro vezes o peso do
corpo num degrau e cerca de doze numa aterrissagem; os \omterm{def:g10:muscles-and-joints:joint}{ligamentos}, que
essas cargas nunca põem à prova, foram o tecido que a torção encontrou.
\end{solution}
